
Une façon
équivalente de décrire ce dernier est la suivante : l'hypothèse que \(X\) est régulier implique que ces anneaux locaux sont factoriels (par le théorème de Auslander--Buchsbaum EGA IV 21.11.1), donc que tout faisceau inversible \(\underline{L}_t\) sur \(X_t\) peut se prolonger en un faisceau inversible \(\underline{L}\) sur \(X\) (EGA IV 21.6.11). Ceci posé, on vérifie aussitôt que l'image par (7.8.1) de \(\cl(\underline{L}_t)\in\Pic(X_t)\) est \(\cl(\underline{L}_s)\in\Pic(X_s)\), où on pose évidemment \(\underline{L}_s=\underline{L}|X_s\).

En fait, n'utilisant que l'hypothèse que les anneaux locaux de \(X\) sont factoriels et que \(X_s\) est localement intègre (à l'exclusion d'hypothèses de régularité et de séparation), on montre que \(\cl(\underline{L}_s)\) ne dépend pas du prolongement choisi \(\underline{L}\) de \(\underline{L}_t\) en un Module inversible sur \(X\), de sorte qu'on obtient ainsi directement un homomorphisme (7.8.1). Pour vérifier l'indépendance annoncée, on est ramené à prouver que si \(\underline{L}_t\simeq\underline{O}_{X_t}\), alors \(\underline{L}_s\simeq\underline{O}_{X_s}\). Or l'hypothèse signifie que \(\underline{L}\simeq\underline{O}_X(D)\), \(D\) un diviseur sur \(X\) de support \(\subset X_s\) (cf. EGA IV 21.2.11), et on se réduit aussitôt au cas \(X_s\) \emph{intègre}, donc \(D=nX_s\), mais alors \(D=\operatorname{div}(g^n)\) (\(g\) une uniformisante de \(A\)), donc \(\underline{O}_X(D)\simeq\underline{O}_X\) et à fortiori \(\underline{L}_s\simeq\underline{O}_{X_s}\).

\subsection*{7.9. Spécialisation et équivalence numérique.}
Supposons \(f\) propre, et \(X_s\) et \(X_t\) réguliers, et considérons les deux applications de spécialisation (7.4.1) et (7.7.4) :
\begin{equation}\tag{7.9.1}
\begin{gathered}
  \sigma_\bullet:\Gr_\bullet(X_t)\longrightarrow\Gr_\bullet(X_s)\\
  \sigma^\bullet:\Gr^\bullet(X_t)\longrightarrow\Gr^\bullet(X_s).
\end{gathered}
\end{equation}
Du fait que l'homomorphisme
\begin{equation}\tag{7.9.2}
  \sigma:K^\bullet(X_t)\longrightarrow K^\bullet(X_s)
\end{equation}
